\begin{center}\begin{tikzpicture}[font=\small,>=stealth,line join=round]
\draw[gray!40,rounded corners=3pt] (0,0) rectangle (14.2,8.3);\node at (7.1,7.700000000000001) {对数定点是(1,0)，不是指数函数的(0,1)};
\begin{scope}[shift={(1.882608695652174,4.029310344827586)},x=1.3043478260869565cm,y=0.793103448275862cm]\draw[->,gray] (-0.6,0)--(4,0) node[right]{$x$};\draw[->,gray] (0,-3)--(0,2.8) node[above]{$y$};\draw[blue!75!black,thick,solid,domain=0.16:3.7,samples=180] plot(\x,{ln(\x)/ln(2)});\fill[orange!85!black] (1,0) circle(2.2pt);\draw[red!70!black,thick] (-.12,.88)--(.12,1.12) (-.12,1.12)--(.12,.88);\draw[gray!35,densely dotted] (1,0)--(1,-3);\node[below=3pt] at (1,-3){$1$};\draw[gray!35,densely dotted] (2,0)--(2,-3);\node[below=3pt] at (2,-3){$2$};\end{scope}\node[anchor=west,align=left,text width=5.0cm] at (8.4,4.0) {正确：$\log_a1=0$\\[10pt]定点$(1,0)$\\[10pt]错误候选$(0,1)$在$y$轴上\\[10pt]真数为0，对数无意义};
\node[text width=12.6cm,align=center] at (7.1,.65) {用“真数为1”找定点；不要把指数函数的共同点直接照搬过来。};
\end{tikzpicture}\end{center}
